\documentclass[12pt]{article}
\usepackage{amsmath}

\usepackage{xcolor}

\usepackage{listings}
\usepackage[most]{tcolorbox}
\usepackage[colorlinks, allcolors=blue]{hyperref}

% style stolen from haskell notes (non avrei ricreato appositamente altrimenti...)
\definecolor{mzbg}{HTML}{F8F9FA}
\definecolor{mzpurple}{HTML}{5E5086}
\definecolor{mzgreen}{HTML}{118C4F}
\definecolor{mzred}{HTML}{C0392B}
\definecolor{mzblue}{HTML}{2E86C1} 
\definecolor{mznumber}{HTML}{ADB5BD}
\definecolor{mztitlebg}{HTML}{2C3E50}

\lstdefinelanguage{MiniZinc}{
    morekeywords={include, var, par, array, of, int, float, set, bool, constraint, solve, satisfy, minimize, maximize, forall, exists, where, if, then, else, endif},
    morekeywords=[2]{alldifferent, abs, sum, count, min, max, card}, 
    sensitive=true,
    morecomment=[l]{\%},
    morecomment=[s]{/*}{*/},
    morestring=[b]",
}

\lstdefinestyle{mznstyle}{
    language=MiniZinc,
    commentstyle=\color{mzgreen}\itshape,
    keywordstyle=\color{mzpurple}\bfseries,
    keywordstyle=[2]\color{mzblue},
    numberstyle=\tiny\color{mznumber},
    stringstyle=\color{mzred},
    basicstyle=\ttfamily\footnotesize,
    breaklines=true,
    showstringspaces=false,
    tabsize=2,
    numbers=none
}

\newtcblisting{mznbox}[1]{
    enhanced,
    skin=bicolor,
    colback=mzbg,
    colbacktitle=mztitlebg,
    coltitle=white,
    colframe=mztitlebg,
    boxrule=0.5pt,
    arc=4pt,
    drop shadow=black!10!white,
    fonttitle=\bfseries\sffamily,
    title=#1,
    listing only,
    listing options={style=mznstyle, numbers=left, numbersep=10pt},
    left=15pt,
    top=4pt,
    bottom=4pt,
    breakable
}

\title{E.A.5.7 - CSP: Cards 2}
\author{}
\date{}

\begin{document}
\maketitle
\vspace{-5em}
\section{Modellazione}

\subsection*{variabili}
        \begin{itemize}
            \item $X =\{ X_i \mid 1\leq i \leq N\}$ - ogni variabile rappresenta il valore della carta nella posizione $i$
            \item $P = \{P_i \mid 1 \leq i \leq 3 \}$ - rappresentano le posizioni dei ``pivot''
        \end{itemize}
        \subsection*{domini}
        \begin{itemize}
            \item $D_{X_i} =\{ 1, \dots, M\}$ per ogni $X_i$
            \item $D_{P_1} = \{ 2, \dots, N - 1 - D \}$ - deve esserci spazio per il primo tratto crescente e per la distanza $D$
            \item $D_{P_2} = \{ 3, \dots, N - 2 \}$ - deve trovarsi tra $P_1$ e $P_3$
            \item $D_{P_3} = \{ 4, \dots, N - 1 \}$ - deve esserci spazio per l'ultimo tratto decrescente
        \end{itemize}
        (potremmo anche scrivere che $D_P = \{1, \dots, N\}$ e far gestire il resto dai vincoli, ma in questo modo restringiamo i domini prima di iniziare la computazione)
    \subsection*{vincoli}
        \begin{itemize}
                \item ordinamento delle sequenze:
            \begin{itemize}
            \item $X_i < X_{i+1}$ per ogni $1 \leq i < P_1$
            \item $X_i > X_{i+1}$ per ogni $P_1 \leq i < P_2$ 
            \item $X_i < X_{i+1}$ per ogni $P_2 \leq i < P_3$
            \item $X_i > X_{i+1}$ per ogni $P_3 \leq i < N$
            \end{itemize}
                \item vincoli sui pivot:
            \begin{itemize}
            \item $P_1 < P_2 < P_3$ - fissiamo un ordinamento
                
            \item $P_3 = P_1 + D$ - distanza tra i due picchi
            \end{itemize}

        \item le carte utilizzate sono quelle della collezione:
            \begin{itemize}
                \item $[X] = Cards$ 
                    \subitem dove $[X]$ è il multiset contenente gli elementi di $X$
                \item (alternativamente esprimibile come:

                    \subitem per ogni valore $v \in \{1, \dots, M\}$

                    $\quad |\{i \in \{1, \dots, N\} \mid X_i = v\}| = |\{j \in \{1, \dots, N\} \mid Cards_j = v\}|$
                    \subitem ovvero, per ogni possibile valore di una carta, il numero di volte in cui quel valore appare in $X$ è uguale al numero di volte in cui appare in $Cards$)
            \end{itemize}
        \end{itemize}

\section{Istanziazione}
$$(Cards, N, M, D) = (\{1, 1, 2, 2, 3, 3, 4\}, 7, 4, 4)$$

\subsection*{variabili e domini}
\begin{itemize}
    \item $X = \{X_1, X_2, X_3, X_4, X_5, X_6, X_7\}$ con $D_{X_i} = \{1, 2, 3, 4\}$ per ogni $i \in \{1, \dots, 7\}$.
    \item $P = \{P_1, P_2, P_3\}$ con:
    \begin{itemize}
        \item $D_{P_1} = \{2\}$ ($N - 1 - D = 7 - 1 - 4 = 2$)
        \item $D_{P_2} = \{3, 4, 5\}$
        \item $D_{P_3} = \{6\}$ ($P_3 = P_1 + D = 2 + 4 = 6$)
    \end{itemize}
\end{itemize}

\subsection*{vincoli}
\begin{itemize}
\item \textbf{ordinamento}
    \begin{enumerate}
        \item $X_1 < X_{2}$         
        \item $X_i > X_{i+1}$ per $2 \leq i < P_2$
        \item $X_i < X_{i+1}$ per $P_2 \leq i < 6$
        \item $X_6 > X_{7}$
    \end{enumerate}

    \item \textbf{distanza e pivot}
    \begin{itemize}
        \item $P_3 = P_1 + 4$
        \item $P_1 < P_2 < P_3$, ovvero $2 < P_2 < 6$
    \end{itemize}

    \item \textbf{carte}
        \begin{itemize}
            \item il multiset $[X_1, X_2, X_3, X_4, X_5, X_6, X_7]$ deve essere uguale a $\{1, 1, 2, 2, 3, 3, 4\}$.
        \item formalmente:
        \begin{itemize}
            \item $|\{i \mid X_i = 1\}| = 2$
            \item $|\{i \mid X_i = 2\}| = 2$
            \item $|\{i \mid X_i = 3\}| = 2$
            \item $|\{i \mid X_i = 4\}| = 1$
        \end{itemize}
    \end{itemize}
\end{itemize}

\section{Codifica in MiniZinc}

\begin{mznbox}{}
include "globals.mzn"; % per global_cardinality

int: M;
int: N;
int: D;
array[1..N] of int: C;

array[1..N] of var 1..M: X;

var 2..N-1-D: P1; % primo picco
var 3..N-2: P2; % minimo locale 
var 4..N-1: P3; % secondo picco

array[1..M] of int: counts = [count(C, v) | v in 1..M];
constraint global_cardinality(X, [v | v in 1..M], counts);

% alternativamente:
% constraint sorted(C) == sorted (X);

constraint P1 < P2;
constraint P2 < P3;
constraint P3 = P1 + D;

constraint forall(i in 1..P1-1) (X[i] < X[i+1]);
constraint forall(i in P1..P2-1) (X[i] > X[i+1]);
constraint forall(i in P2..P3-1) (X[i] < X[i+1]);
constraint forall(i in P3..N-1) (X[i] > X[i+1]);

solve satisfy;
\end{mznbox}

\end{document}
